\section{引言}

人类视觉系统的价值，并不只体现在“认出了什么”，更体现在“构成了怎样的世界”。从视网膜上的二维投影到对物体位置、朝向、距离、遮挡、可达性与未来状态的判断，视觉的深层功能一直是把局部、瞬时、噪声化的感受，转化为一个足够稳定的外部场景模型。正因为如此，在心理物理学、系统神经科学与 psychometrics 中，研究者很少满足于单个刺激上的正确率；他们更关心观察者在系统性探测、重复测量、条件扰动和跨任务比较下，是否表现出一致的结构、稳定的表征与可解释的失真模式 \cite{kriegeskorte2008rsa,yamins2016dicarlo,xu2025bsa,zhang2025comfort}。如果我们真的希望理解 VLM “看见了什么”，那么就不该只把它当作一个会答题的黑箱，而应该像研究一个观察者那样，去追问：它的回答背后是否对应着一个足够一致、足够可重建的世界。

现有 VLM 评测与这一目标之间存在明显落差。大量工作把空间理解压缩成单题问答、选择题正确率或局部关系判断：模型说对了“左边”“更近”“前方”这样的答案，就被视为具备某种空间理解能力。然而，答对局部问题，并不等于形成了全局场景；多个单独正确的回答，也完全可能彼此不兼容。一个模型也许能在单图条件下猜对一部分距离比较，却无法把这些判断拼接成一致的布局；它也可能在语言模板的帮助下给出看似合理的方位答案，但这些答案一旦进入统一坐标系，就会立即暴露出冲突、塌缩和不可实现性。换句话说，当前许多 benchmark 测到的，其实更像是“局部关系回答能力”，而不是“可被重建的空间信念”。

本文的核心动机来自序数嵌入与场景重建的一个基本事实：当偏序关系足够充分且足够准确时，场景几何可以在相似变换意义下被近似恢复 \cite{agarwal2007gnmds,ma2018robustordinal,dokmanic2015edm}。这意味着，如果我们不再满足于向 VLM 提几个零散问题，而是系统性地询问大量相对距离与相对方位关系，那么模型的回答集合本身就构成了一种行为学意义上的“断层扫描”。在这个视角下，QRR 与 TRR 不只是两个问答任务，而是两类读取模型外显场景信念的探针：QRR 通过距离偏序检验模型是否保留了物体间的相对近远结构；TRR 通过方向关系检验模型是否维持了统一的参照系与方位几何。若这些回答能够支撑一个与真实场景接近的重建结果，我们才有理由说模型确实形成了某种场景级表征；反之，如果这些回答只能生成一个高度冲突、严重失真或完全不可实现的布局，那么模型即使在若干题目上答对了，也更可能只是依赖语言先验、模板匹配或彼此独立的局部启发式。

这也使 Ordinary-Bench 与两类看似相近、但目标不同的研究形成了清楚区分。第一类是 world model 工作：从经典的 latent dynamics 建模，到近期用于导航、驾驶和 4D 场景预测的生成式模型，它们的核心目标通常是预测未来观测、生成视频、BEV 或 occupancy，并服务于规划与控制 \cite{ha2018worldmodels,wang2024drivewm,bar2025navigationwm,ni2025maskgwm,liao2025i2world,zhou2025hermes}。第二类是 visual CoT / tool-use 路线：模型通过外部视觉工具、代码执行、视觉草图板、树搜索或更长的 test-time reasoning 来把题做对 \cite{suris2023vipergpt,liu2023llavaplus,gui2024vpd,hu2024visualsketchpad,wang2025visuothink,lin2025multimodalcot}。这两条路线都很重要，但它们优化的主要是生成、规划或解题流程，而不是回答本身是否已经定义了一个对象级、可反演、可矢量化的场景结构。Ordinary-Bench 要测的，正是这一更基础的问题。

基于这一思路，我们提出 \textbf{Ordinary-Bench}。它不是再增加一个“空间题库”，而是把 VLM 当作一个需要被行为反演的观察者：我们在受控 3D 渲染环境中生成结构明确的场景，以 QRR（四元相对关系）和 TRR（三元时钟方向关系）进行密集探测，再将模型输出转写为偏序约束与方向约束，检验这些回答是否足以支持一个一致的场景重建。使用受控合成环境而不直接依赖自然图像，是一个刻意的选择。对于“模型究竟构成了什么场景”这一问题，精确几何真值、可控复杂度和可重复扰动，比生态真实性更重要；只有在场景结构本身无歧义的条件下，我们才能严格地区分视觉贡献、语言先验和结构性失真 \cite{johnson2017clevr,yi2020clevrer}。

当前已经完成的实验结果进一步说明了为什么需要这种更严格的评估框架。基于目前整理出的 540 个场景与 200000+ QA 对，我们观察到：在单图条件下，模型在 QRR 上还能达到大约 70\% 的准确率，看起来似乎具备了某种距离比较能力；但一旦输入切换到多视角，QRR 表现反而退化到接近随机猜测，表明模型并没有把多视角信息整合为更稳定的空间表征，而更像是在额外视角的干扰下丧失了原本脆弱的局部线索。TRR 的情况更差，几乎始终停留在随机猜测水平。这意味着当前 VLM 对“谁更近”可能还有有限的、脆弱的判断能力，但对“统一参照系中的相对方向”几乎没有形成可靠表征。更值得注意的是，基于 540 个场景、200000+ QA 对进行微调或强化学习，并没有带来显著的提高和优化。这一负结果具有重要含义：仅仅把更多关系问答灌输给模型，并不会自动产生几何一致的场景信念；如果评测目标仍停留在局部正确率上，这种结构性失败很容易被掩盖。

从更大的应用角度看，这一问题并不只是“空间题有没有答对”。如果一个 VLM 真的能在偏序关系层面形成可重建、可想象、可对齐的场景信念，那么它实际上已经具备了一种与具体数值尺度弱耦合的场景矢量化能力。偏序重建并不要求模型先恢复每一条精确欧氏距离；只要关系结构稳定存在，我们就可以在后处理中再引入真值尺度、经验尺度或任务相关尺度，把相对布局对齐为可用于定位、检索、导航、交互甚至场景编辑的结构化表示。也就是说，Ordinary-Bench 评估的并不是一个狭义的 QA 技能，而是在测试 VLM 是否已经触碰到“从视觉中恢复可操作场景结构”这一更本质的能力。

进一步地，Ordinary-Bench 不应只停留在静态场景。既然真实视觉的目标是维持一个随时间更新的场景世界，那么更合理的实验设计应该同时覆盖：单图下的局部关系感知，多视角下的统一场景整合，以及动态场景中关系随时间变化后的持续一致性。也就是说，未来实验不只应问“当前谁更近、谁在谁的左边”，还应问“何时发生关系翻转、何时出现交叉或遮挡、模型在时间推进中是否仍维持同一个场景解释”。这一点也与近期视频与 world-model 评测中对 motion-level、temporal consistency 和 future-state prediction 的关注形成呼应 \cite{hong2025motionbench,fu2024videomme,zhou2024mlvu,liu2024tempcompass,gao2025wmabench}。

本文的贡献可以概括为四点。第一，我们提出一个受心理物理学启发的 VLM 空间评估框架，用受控 3D 场景和密集 QRR/TRR 探测，把模型从“答题器”重新定义为“可被行为反演的观察者”。第二，我们将序数嵌入中的可重建性思想引入 VLM 评测：通过足够多的偏序问答，读取并检验模型的外显场景信念，而不再只报告单题正确率。第三，我们给出一组当前已经相当清晰的负面证据：单图下有限的 QRR 成功并不意味着模型形成了稳定空间表征；多视角并未提升而是压垮了这种能力，TRR 长期接近随机，且微调与强化学习并未实质改变这一局面。第四，我们进一步提出一个从静态到动态的实验设计方向，把多视角整合、时序关系变化、场景级可实现性与工具增强对照纳入统一 protocol。综合来看，本文试图把问题从“VLM 会不会回答空间问题”推进到一个更根本也更严格的问题：\emph{VLM 是否真的在视觉输入上构造了一个足够一致、足够可复原的世界？}
